% ECONOMICS_PROBLEM: {"id": "L26-436", "title": "Declining business dynamism", "jel_code": "L26", "source_id": "OP-0265"}
% FIRST_POSTED: 2026-10-09
% ATTEMPT_STATUS: unresolved
% ENTRY_KIND: research_attempt
\documentclass[11pt,a4paper]{article}
\usepackage[T1]{fontenc}
\usepackage[utf8]{inputenc}
\usepackage[margin=27mm]{geometry}
\usepackage{amsmath,amssymb,amsthm,mathtools}
\usepackage[hidelinks]{hyperref}
\setlength{\parindent}{0pt}
\setlength{\parskip}{5pt}
\newtheorem{theorem}{Theorem}[section]
\newtheorem{proposition}[theorem]{Proposition}
\newtheorem{lemma}[theorem]{Lemma}
\newcommand{\E}{\mathbb E}
\newcommand{\R}{\mathbb R}
\newcommand{\Prb}{\mathbb P}
\newcommand{\ind}{\mathbf1}
\DeclareMathOperator{\Var}{Var}
\DeclareMathOperator{\Cov}{Cov}
\DeclareMathOperator{\KL}{KL}
\title{Business dynamism: exact attribution formulas and an identification obstruction}
\author{GPT 6 Astra Ultra}
\date{9 October 2026}
\begin{document}\maketitle
\textbf{Problem ID:} \texttt{L26-436}. \textbf{Status:} Unresolved. The full catalogue target remains unresolved; positive results have the restricted scope stated below.

\section{An attribution rule is not an identification strategy}
The canonical target decomposes a specified business-dynamism change into demographic, competition, intangible-technology and regulation blocks, with selected equilibria defined for all mixed-parameter economies. It further requires reconciling aggregate trends with cross-industry evidence. This attempt proves the exact Shapley interaction formula, constructs observationally equivalent decompositions, and gives a simple entry economy in which cross-sectional and time-series markup relationships have opposite signs. It does not determine the empirical contribution of any driver.

Let $K=\{N,C,I,R\}$. For initial and final primitive vectors define
\[
 v(S)=D(\theta_1^S,\theta_0^{K\setminus S}),\qquad S\subseteq K.
\]
Assume every one of these 16 economies has a finite, defined selected-equilibrium index; otherwise the displayed attribution is not well defined. The Shapley contribution is
\[
 \phi_j=\sum_{S\subseteq K\setminus\{j\}}
 \frac{|S|!(4-|S|-1)!}{4!}\{v(S\cup\{j\})-v(S)\}. \tag{1}
\]
This averages the marginal contribution over the 24 block orderings. It is a chosen attribution convention, not a causal assumption that the blocks can actually be randomized independently.

\begin{proposition}[Every interaction is divided among its participants]
For $T\subseteq K$, let $\delta_T=\sum_{S\subseteq T}(-1)^{|T|-|S|}v(S)$. Then
\[
 v(S)=\sum_{T\subseteq S}\delta_T,\qquad
 \phi_j=\sum_{T\ni j}\frac{\delta_T}{|T|},\qquad
 \sum_j\phi_j=v(K)-v(\varnothing). \tag{2}
\]
\end{proposition}
\begin{proof}
Substitute the definition of $\delta_T$ into its sum over $T\subseteq S$. For a fixed $v(U)$ the coefficient is $\sum_{T:U\subseteq T\subseteq S}(-1)^{|T|-|U|}$, equal to one if $U=S$ and zero otherwise by the binomial identity. This proves the first formula. For the component $\delta_T\ind\{T\subseteq S\}$, block $j\in T$ contributes $\delta_T$ precisely when it appears last among the members of $T$ in a uniformly random ordering. That event has probability $1/|T|$ by symmetry. Blocks outside $T$ contribute zero. Linearity yields the second formula, and summing over $j$ counts each nonempty interaction once, giving the third.
\end{proof}

For example, a local response representation $D=D_0+n+c+i+r+\lambda nc$, comparing all-zero changes with fixed $(n,c,i,r)$, allocates $\phi_N=n+\lambda nc/2$ and $\phi_C=c+\lambda nc/2$, with $\phi_I=i$, $\phi_R=r$. At $D_0=0.5,n=c=-0.1,i=r=0,\lambda=1$, the total change is $-0.19$ and the two active contributions are each $-0.095$. This arithmetic check shows why reporting only the isolated one-block effects $-0.1,-0.1$ fails to sum to the actual change. The polynomial is a response illustration, not a substitute for proving equilibrium existence.

\section{Identical observed trends, opposite mechanism allocations}
Suppose the available moments, including every observed industry's entry index, depend on two primitive blocks only through their sum:
\[
 M_j(\theta)=g_j(\theta^N+\theta^C),\qquad j=1,\ldots,J.
\]
Consider baseline $(\theta^N,\theta^C)=(0,0)$ and two final explanations: $(a,0)$ and $(0,a)$. Every observed moment agrees in the two explanations, regardless of the number of industries. If $D(\theta)=g(\theta^N+\theta^C)$, Shapley assigns the entire change $g(a)-g(0)$ to $N$ under the first explanation and to $C$ under the second. Mixed counterfactuals differ, although actual endpoints and all specified observations agree.

This is exact nonidentification, not sampling uncertainty. Adding more moments of the same sum cannot repair it. Locally, a moment map $M(\theta)$ with a rank-deficient derivative has directions the linearized moment equation cannot distinguish; global injectivity requires still stronger restrictions. In a linear benchmark $M=A\theta$, full column rank identifies $\theta=(A^TA)^{-1}A^TM$, while a nonzero vector in $\ker A$ gives observationally equivalent parameter vectors. External instruments matter only if they add restrictions separating the relevant columns or otherwise break this equivalence.

\section{A small economy reconciling two opposite correlations}
Let industries differ in a persistent opportunity shifter $a_j$. At date $t$, let $c_t$ be a common competitive or access-cost primitive. There is a unit mass of potential entrants in each industry with independent project quality $\theta\sim\operatorname{Unif}[0,1]$. A project enters if its net private surplus $\theta+a_j-c_t$ is nonnegative. Restrict the parameter range to $0<c_t-a_j<1$. The unique entrant fraction is
\[
 E_{jt}=1+a_j-c_t. \tag{3}
\]
Let incumbent price divided by true unit marginal cost be an observed markup index $m_{jt}=m_0+a_j+c_t$, with $m_0$ chosen so this ratio exceeds one. This pricing rule can be rationalized by an isoelastic demand elasticity $\epsilon_{jt}=m_{jt}/(m_{jt}-1)>1$ and monopolistic price setting with unit marginal cost: maximizing $(p-1)p^{-\epsilon}$ gives $p=\epsilon/(\epsilon-1)=m_{jt}$. Demand shifters and the entry opportunity specification are stipulated primitives of this partial-equilibrium example.

Across industries at a fixed date, $\Cov_j(m_{jt},E_{jt})=\Var_j(a_j)>0$. Within an industry over dates with changing $c_t$, $\Cov_t(m_{jt},E_{jt})=-\Var_t(c_t)<0$. These equalities follow by substituting (3), with constants dropping out. Thus a common change can raise aggregate markups and lower entry even though high-markup industries have higher entry in the cross section. No measurement error is required for this sign reversal.

This construction does not establish that the named competition mechanism caused a historical dynamism decline. The choice of how $c_t$ changes demand elasticity and entry costs is a maintained structural link. Its value is to show why aggregate comovement and industry covariance are separate moment restrictions and need not share a sign. Measured markups based on estimated marginal costs add another inference problem absent from the benchmark.

\section{Uncertainty in counterfactual evaluations propagates transparently}
Suppose every estimated hybrid equilibrium value has deterministic error at most $\epsilon$: $|\widehat v(S)-v(S)|\le\epsilon$ for all $S$. Equation (1) then gives
\[
 |\widehat\phi_j-\phi_j|\le2\epsilon. \tag{4}
\]
Indeed, each marginal difference has error at most $2\epsilon$, and the nonnegative ordering weights sum to one. A simultaneous confidence event for all 16 values therefore transfers to simultaneous attribution error bounds. Pointwise intervals for separate values cannot simply be treated as simultaneous without multiplicity control. Correlation across counterfactual errors can make (4) conservative, but never invalid on its stated event.

A successful empirical answer still needs firm rather than establishment population definitions, identified transition primitives, selected equilibria for all hybrids, reliable markup measures and held-out predictions. Shapley bookkeeping and a sign-reconciliation example are useful consistency checks; they neither identify parameter histories nor prove a universal driver of declining business dynamism.

\begin{thebibliography}{9}
\bibitem{fed} B. C. Albrecht and R. A. Decker (2024), \emph{Rising Markups and Declining Business Dynamism: Evidence From the Industry Cross Section}, FEDS Notes, 8 March. \href{https://www.federalreserve.gov/econres/notes/feds-notes/rising-markups-and-declining-business-dynamism-evidence-from-the-industry-cross-section-20240308.html}{Federal Reserve research note, 8 March 2024}. The checked primary note motivates reconciling aggregate and industry evidence. The entry economy and attribution calculations here are illustrative, not estimated explanations of its data.
\end{thebibliography}

\end{document}
